\chapter{Asia-Specific Equity Strategies}\label{ch:s1:asia-specific-equity-strategies}

A Shanghai stock reaches its ten per cent limit up at 10:05 and stays there all day; the buyers who could not trade are still there tomorrow, and it
opens higher. Seasholes and Wu found that such events draw individual investors who did not own the stock, that prices rise and then revert over the
following week, and that smart traders who bought on the limit day and sold the next earned an average of 1.16\% a day. Chen, Gao, He, Jiang and Xiong
found large investors doing the same, and pushing prices to the limit to do it. On this chapter's synthetic limit market, a buyer at every limit-up
close would gain 4.5\% by the next open; a buyer who joins the queue at the limit gains 2.6\%, because the queue fills when the lock is weakest; and
the price gives back 2.8\% over the following week. The build is \texttt{firm.limitmkt}.

\section{Price limits and T+1}

Book~1, chapter~12 described the rules: mainland main-board shares may move 10\% a day, those on the growth boards 20\%, and shares bought cannot be
sold the same day (its dated box has the current widths). It also showed what limits do to the data: a limit splits a large move over several days, so
observed returns after a limit day continue, and a signal that sees the continuation in closing prices has found a return that cannot be bought at
those prices. This chapter is about the trades that can.

\begin{definition}[Limit-hit continuation]\label{def:s1:asia-specific-equity-strategies:continuation}
\emph{Limit-hit continuation}\index{limit-hit continuation} is the tendency of a stock that closes locked at its daily limit to move further in the
same direction at the next open, because the part of the move the limit held back, and the orders queued at the limit, carry over to the next
day.
\end{definition}

\texttt{firm.limitmkt} prints a limit market from fair returns (\cref{lst:s1:asia-specific-equity-strategies:limits}). It takes the synthetic
market's ten years on 1,000 names, scales their returns by 1.5 to a daily volatility of 3.1\%, and moves each close toward the fair value plus an
attention premium, but never beyond the limit. A close at the upper limit draws attention buyers the next morning, adding 3\% to the target, a premium
that decays by 0.6 a day. The part of a move beyond the limit waits for the next open. And a buy order joining the queue at a locked limit fills
with a probability that falls with how far beyond the limit the stock's value is: a lock that barely holds fills, one with a large backlog does not.
That last rule is the queue's adverse selection. The buyers most likely to be filled are those who least wanted to be.

\begin{definition}[Magnet effect]\label{def:s1:asia-specific-equity-strategies:magnet}
The \emph{magnet effect}\index{magnet effect} is the tendency of prices to accelerate toward a daily price limit as they approach it, more strongly
toward the upper limit than the lower, as traders hurry to trade before the limit stops them or push the price to lock.
\end{definition}

Cho, Russell, Tiao and Tsay documented the \omterm{def:s1:asia-specific-equity-strategies:magnet}{magnet effect} in Taiwan; Chen and co-authors explained it with large investors who push prices to the upper
limit and sell the next day. The layer plants their mechanism: half the targets within 2\% below the upper limit are pushed to lock, closing above
their fair value. \Cref{fig:s1:asia-specific-equity-strategies:magnet} shows the footprint: with the magnet, days that would close 7.5\% to 9.5\% up
are half as frequent and limit-up closes rise from 0.87\% of stock-days to 1.29\%.

\begin{omfigure}
\begin{tikzpicture}
\begin{axis}[width=10cm,height=5.2cm,ybar=0pt,bar width=6pt,xlabel={\small close-to-close log move (\%)},ylabel={\small per 10,000 stock-days},
  tick label style={font=\footnotesize},xmin=4.9,xmax=9.6,ymin=0,ymax=110,grid=major,grid style={gray!25},legend pos=north east,
  legend style={font=\scriptsize},area legend,table/col sep=comma]
\addplot[fill=omExo!50,draw=omExo] table[x=bin_pct,y=none]{figdata/strategies-1/18-asia-specific-equity-strategies/magnet.csv};
\addplot[fill=omDef!70,draw=omDef] table[x=bin_pct,y=magnet]{figdata/strategies-1/18-asia-specific-equity-strategies/magnet.csv};
\legend{no magnet,with the magnet}
\end{axis}
\end{tikzpicture}

\omcaption{Daily moves just below a 10\% limit on the synthetic limit market, in half-point bins, with and without the magnet (half of the targets within
2\% of the limit pushed to lock). The pushed days leave the bins from 7.5\% and become limit-up closes: 129 per 10,000 stock-days with the magnet, 87
without. Data: \texttt{s1\_asia.band}.}\label{fig:s1:asia-specific-equity-strategies:magnet}
\end{omfigure}

The next table follows every limit-up close for a week (\cref{lst:s1:asia-specific-equity-strategies:ups}). The trade is the one the rules allow:
buy at the limit on the lock day, sell at the next open. It is also what the T+1 rule makes necessary, since the shares cannot be sold the same day.

\begin{center}
\small
\begin{tabular}{@{}lccc@{}}
\toprule
limit-up closes at a 10\% limit & all & value beyond the limit & pushed\\
\midrule
share of the closes & 100\% & 67\% & 33\%\\
return to the next open, every close & 4.51\% & 5.90\% & 1.67\%\\
queue fill probability & 57\% & 36\% & 100\%\\
return to the next open, weighted by fills & 2.64\% & 3.96\% & 1.67\%\\
next open to the close five days later & $-2.81\%$ & $-2.77\%$ & $-2.88\%$\\
\bottomrule
\end{tabular}
\end{center}

Two effects make the overnight gain. The backlog is the part of the move held back by the limit; without attention or magnet it gives 3.42\% to the next
open (1.21\% after fills). The attention buyers add most of the rest, and take it back over the week. The pushed locks earn little overnight, because they
closed above their value, but they still draw the attention buyers. Fills cut the gain from 4.51\% to 2.64\%: the buyer who could not be filled in the
strongest locks is the rule, not the exception. At a round trip of 15 basis points (the sell-side stamp duty of the dated box plus assumed commissions
and spread) the filled trade still clears 2.49\% on average, about twice the smart traders' 1.16\% in Seasholes and Wu, because the model's backlog and
attention are clean and every queue position is equal.

\begin{omfigure}
\begin{tikzpicture}
\begin{axis}[width=10cm,height=5.4cm,ybar,bar width=10pt,xlabel={\small daily price limit (\%)},ylabel={\small mean return (\%)},
  tick label style={font=\footnotesize},symbolic x coords={5,10,20},xtick=data,ymin=-6,ymax=8,enlarge x limits=0.25,grid=major,
  grid style={gray!25},legend columns=3,legend style={font=\scriptsize,at={(0.5,-0.3)},anchor=north,draw=none,fill=none,
  /tikz/every even column/.append style={column sep=8pt}},area legend,table/col sep=comma]
\addplot[fill=omThm!60,draw=omThm] table[x=limit_pct,y=open_pct]{figdata/strategies-1/18-asia-specific-equity-strategies/widths.csv};
\addplot[fill=omDef!60,draw=omDef] table[x=limit_pct,y=filled_pct]{figdata/strategies-1/18-asia-specific-equity-strategies/widths.csv};
\addplot[fill=omExo!50,draw=omExo] table[x=limit_pct,y=week_pct]{figdata/strategies-1/18-asia-specific-equity-strategies/widths.csv};
\legend{close to next open,weighted by queue fills,next open to five days later}
\end{axis}
\end{tikzpicture}

\omcaption{Limit-up closes on the synthetic limit market by the width of the daily limit: the mean return from the close to the next open for every close
and weighted by the probability that a queued buy order fills, and the mean return over the following week. Data:
\texttt{s1\_asia.limit\_ups}.}\label{fig:s1:asia-specific-equity-strategies:widths}
\end{omfigure}

The width decides how often and how hard (\cref{fig:s1:asia-specific-equity-strategies:widths}). A 5\% limit locks on 11.24\% of stock-days (the
synthetic stocks are volatile), a 10\% limit on 1.29\%, a 20\% limit on 0.10\%. The wider the limit, the more extreme the move behind each lock: the
next-open gain rises from 2.98\% to 4.51\% and 6.58\%, the filled gain from 2.12\% to 2.64\% and 3.20\%, and the week's give-back from 1.76\% to 2.81\% and
5.27\%. The 20\% column rests on 2,325 events and its week-after figure is noisy; the next-open means have standard errors of 0.02 points at 10\% and
0.1 at 20\%.

\begin{dated}{2026-09}{Mainland costs and Connect data}\label{dat:s1:asia-specific-equity-strategies:connect}
From 28 August 2023 the stamp duty on trading Shanghai and Shenzhen shares was halved, from 0.1\% to 0.05\% of the transaction amount, paid by the
seller. On 12 April 2024 the Hong Kong, Shanghai and Shenzhen exchanges announced changes to northbound Stock Connect data, in two phases about one
and four months later: real-time buy, sell and total turnover would no longer be available (daily and monthly totals and the ten most active stocks
would be); the remaining daily quota would be shown only when below 30\%; and northbound shareholdings would be published quarterly, on the fifth
northbound trading day after the quarter's end.
\end{dated}

\section{Retail-dominated markets}

Most trading in mainland shares is done by individuals, and they are not one group. Jones, Shi, Zhang and Zhang split Chinese retail accounts by
size: the smallest cannot predict returns, chase the day's moves, fail to process public news and show gambling preferences; the largest predict
returns correctly and trade against the day's moves, though their costs absorb what they gain. The attention buying after limit-up days is the
small accounts' behaviour, and the large investors of Chen and co-authors are on the other side of it.

The week after a limit-up close is where the retail strategy lives. In the synthetic market the attention premium reverses by 2.81\% from the next open
to five days later. Harvesting it takes a short sale, and borrow in these markets can be scarce or, as in Korea's seventeen months without short
selling (Book~1, chapter~12), banned outright. In practice the trade is the other side: a holder sells into the attention at the open, and a book that would have
bought on the continuation stays out.

\section{Connect flows}

\begin{definition}[Northbound flow]\label{def:s1:asia-specific-equity-strategies:northbound}
\emph{Northbound flow}\index{northbound flow} is the net buying of mainland shares by investors trading through Stock Connect from Hong Kong, by
stock or in aggregate; \emph{southbound} flow is the reverse, mainland investors buying Hong Kong shares.
\end{definition}

Northbound investors are mostly institutions, and Bian, Chan, Han and Shi found that their net purchases predict the returns of connected Shanghai
stocks. A strategy that follows them needs their flows by stock, and soon: information in a flow is used up quickly. The dated box records that the
by-stock holdings, once published daily, are now quarterly. The chapter plants flows that know a little of each stock's next five days (a correlation of
0.02 with the return) and trades a weekly long--short book on them, at 15 basis points a round trip.

\begin{center}
\small
\begin{tabular}{@{}lcc@{}}
\toprule
northbound holdings published & daily & quarterly, five days late\\
\midrule
IC with the next five days' return & 0.021 & 0.003\\
Sharpe ratio after costs, weekly book & 2.04 & 0.26\\
\bottomrule
\end{tabular}
\end{center}

With daily data the flow's IC of 0.021 over a thousand names a week gives a Sharpe ratio of 2.0 after costs; with quarterly data five days late,
the same flows are worth almost nothing. The small IC that remains (0.003, a standard error of 0.001) comes from the synthetic market's persistent
drifts, which make last quarter's informed buying weakly informative about next week. The lesson is general: a signal made of someone else's
published trades is worth what the publication schedule leaves of it.

\section{Strategy files}

\begin{strategyfile}{Limit-hit continuation}\label{strat:s1:asia-specific-equity-strategies:limit}
\sfield{Who pays you, and why} Buyers held back by the limit, and attention buyers who arrive the next morning.
\sfield{Instruments and venues} Mainland A-shares at their daily limit; other markets with daily limits.
\sfield{Signal} A limit-up close, preferably with a large backlog of unfilled buy orders and news behind it.
\sfield{Sizing and execution} A buy order in the queue at the limit early, sold at the next open (the T+1 rule allows no sooner).
\sfield{Costs} Stamp duty on the sale, commissions; the queue's adverse selection, which is the main cost.
\sfield{How it dies} Competition for queue priority; regulators' attention to large buyers at the limit.
\sfield{Horizon, capacity, infrastructure} Overnight; capacity limited by what the queue fills; order entry fast enough to be early in the queue.
\sfield{Backtest honestly} Fills from the queue, not from closing prices; a close at the limit is not a price anyone could buy at.
\sfield{Sources} Seasholes and Wu (2007): 1.16\% a day to smart traders; Chen, Gao, He, Jiang and Xiong (2019).
\end{strategyfile}

\begin{strategyfile}{Retail-attention reversal}\label{strat:s1:asia-specific-equity-strategies:retail}
\sfield{Who pays you, and why} Small retail accounts that buy attention-grabbing stocks after limit-up days.
\sfield{Instruments and venues} Stocks after limit-up days; for a short, those that can be borrowed.
\sfield{Signal} A limit-up close yesterday, the open's gap, the retail share of trading.
\sfield{Sizing and execution} Sell holdings or short at the open; hold the week.
\sfield{Costs} Borrow, where available at all; stamp duty.
\sfield{How it dies} Short-sale restrictions; less retail trading.
\sfield{Horizon, capacity, infrastructure} A week; capacity limited by borrow.
\sfield{Backtest honestly} Only stocks that could be borrowed on the day; the open's price, not the previous close.
\sfield{Sources} Seasholes and Wu (2007): mean reversion over the following week; Jones, Shi, Zhang and Zhang (2025).
\end{strategyfile}

\begin{strategyfile}{Northbound-flow signal}\label{strat:s1:asia-specific-equity-strategies:northbound}
\sfield{Who pays you, and why} Investors slower than the institutions trading through the Connect, who learn their information later.
\sfield{Instruments and venues} Connect-eligible mainland shares.
\sfield{Signal} Changes in northbound holdings by stock, or aggregate \omterm{def:s1:asia-specific-equity-strategies:northbound}{northbound flow}.
\sfield{Sizing and execution} A long--short book on the changes, rebalanced as often as the data allow.
\sfield{Costs} Stamp duty and commissions on each rebalance.
\sfield{How it dies} The data: by-stock holdings are now published quarterly.
\sfield{Horizon, capacity, infrastructure} Days to weeks with daily data; the data's schedule sets it.
\sfield{Backtest honestly} Holdings as published on the day they were published, with the schedule in force at the time.
\sfield{Sources} Bian, Chan, Han and Shi (2023); the dated box.
\end{strategyfile}

\section{Tutorial: locked limit up}\label{tut:s1:asia-specific-equity-strategies:tut}

\begin{tutorial}
\textbf{Goal.} Print the synthetic market through daily price limits, follow every limit-up close, and measure what a buyer at the limit, a buyer in
the queue and a seller a week later would have had. \textbf{End state:} the tables and the two figures.
\begin{tutsteps}
\item \textbf{The limit layer}: open, close, locks, the magnet and the queue fills, day by day.
\omcode{code/firm/limitmkt/firm_limitmkt.py}{43}{73}{Printing a limit market from fair returns.}\label{lst:s1:asia-specific-equity-strategies:limits}
\item \textbf{The events}: every limit-up close with a week of listing after it, split into locks with a backlog and pushed locks.
\omcode{code/strategies-1/18-asia-specific-equity-strategies/python/s1_asia.py}{47}{64}{Limit-up closes and what follows them.}\label{lst:s1:asia-specific-equity-strategies:ups}
\item \textbf{Run} \texttt{limit\_ups(limit)} for the three widths and the two controls (\texttt{attention=0}, \texttt{push=0}), \texttt{band(push)},
\texttt{flows(1, 0)} and \texttt{flows(63, 5)}, and \texttt{fig\_asia.py}.
\end{tutsteps}
\textbf{What to change next.} Give queue positions: early buyers fill first; make the attention premium depend on how many stocks lock the same day
(Seasholes and Wu found fewer buyers when many do); add lower-limit locks and the sellers they trap.
\end{tutorial}

\section{Build: a limit market}\label{bld:s1:asia-specific-equity-strategies:limitmkt}

\begin{build}
\bfield{Purpose} Prices of a daily-limit market from fair returns, with attention buying, a magnet and queue fills; a northbound-flow generator and
its publication schedule.
\bfield{Interface} \texttt{LimitConfig(\dots)}, \texttt{apply\_limits(R, cfg, rng)}, \texttt{northbound(alpha, skill, rng)}, \texttt{holdings\_seen(flow,
every, lag)}.
\bfield{Rules} No close beyond the limit from the previous close; the open carries yesterday's backlog and today's premium, within the limit; a listing
starts at its fair value.
\bfield{Acceptance tests} \texttt{code/firm/limitmkt/tests/}: a 15\% move locks at 10\% with the fill probability by hand and opens at the fair value
plus the premium; the magnet pushes a 9\% target; holdings published every four days two days late; the flows' correlation.
\bfield{Stretch} Queue priority; lower-limit locks; a T+1 constraint on intraday round trips in the backtester.
\end{build}

\begin{omsources}
\item T.~Chen, Z.~Gao, J.~He, W.~Jiang and W.~Xiong, ``Daily price limits and destructive market behavior'', \emph{Journal of Econometrics} 208(1),
2019.
\item D.~D.~Cho, J.~Russell, G.~C.~Tiao and R.~Tsay, ``The magnet effect of price limits: evidence from high-frequency data on Taiwan Stock Exchange'',
\emph{Journal of Empirical Finance} 10(1--2), 2003.
\item M.~S.~Seasholes and G.~Wu, ``Predictable behavior, profits, and attention'', \emph{Journal of Empirical Finance} 14(5), 2007.
\item C.~M.~Jones, D.~Shi, X.~Zhang and X.~Zhang, ``Retail trading and return predictability in China'', \emph{Journal of Financial and Quantitative
Analysis} 60(1), 2025 (online 2024).
\item J.~Bian, K.~Chan, B.~Han and D.~Shi, ``Cross-border equity flows and information transmission: evidence from Chinese stock markets'',
\emph{Journal of International Financial Markets, Institutions and Money} 84, 2023.
\end{omsources}

\section{Exercises}

\begin{exercise}[$\star$]\label{exo:s1:asia-specific-equity-strategies:1}
News raises a stock's value by 50\%. How many days does it take to reach it under a 10\% limit, and under a 20\% limit?
\end{exercise}

\begin{exercise}[$\star$]\label{exo:s1:asia-specific-equity-strategies:2}
A stock's value rises 15\% in a day under a 10\% limit. With the layer's fill scale of 2\%, what is the probability that a queued buy order fills?
\end{exercise}

\begin{exercise}[$\star$]\label{exo:s1:asia-specific-equity-strategies:3}
The attention premium is 3\% at the open after a limit-up close and decays by 0.6 a day. How much of it remains four days later, and how much has
reverted?
\end{exercise}

\begin{exercise}[$\star\star$]\label{exo:s1:asia-specific-equity-strategies:4}
Why do the pushed locks have a fill probability of one and a small overnight gain?
\end{exercise}

\begin{exercise}[$\star\star$]\label{exo:s1:asia-specific-equity-strategies:5}
Why does the overnight gain grow with the width of the limit?
\end{exercise}

\begin{exercise}[$\star\star$]\label{exo:s1:asia-specific-equity-strategies:6}
At a 10\% limit, 1.29\% of stock-days are limit-up closes. How many does a market of 1,000 stocks produce a day, and what does the filled trade earn on
each after a 15 basis point round trip?
\end{exercise}

\begin{exercise}[$\star\star\star$]\label{exo:s1:asia-specific-equity-strategies:7}
\emph{Coding.} Run \texttt{limit\_ups(0.10, attention=0.0, push=0.0)}. What is left of the overnight gain, filled and unfilled, and of the week's
reversal? What does that tell you about where the trade's return comes from?
\end{exercise}

\begin{exercise}[$\star\star\star$]\label{exo:s1:asia-specific-equity-strategies:8}
\emph{Find the flaw.} ``Stocks that close limit up gain 4.5\% by the next open on average; buying every limit-up close is a Sharpe ratio of ten.''
\end{exercise}

\section{Problem: Locked Limit Up}

\begin{problem}[{Weekend problem --- a day at the limit}]\label{pb:s1:asia-specific-equity-strategies:1}
The chapter's synthetic limit market and the public record.

\medskip\noindent\textbf{Part I --- The rules.}
\begin{enumerate}
\item Define \omterm{def:s1:asia-specific-equity-strategies:continuation}{limit-hit continuation} and the \omterm{def:s1:asia-specific-equity-strategies:magnet}{magnet effect}.
\item What do daily limits do to observed returns (Book~1, chapter~12)?
\item What does the T+1 rule forbid, and what trade does it leave?
\item Describe the chapter's limit layer.
\end{enumerate}

\medskip\noindent\textbf{Part II --- The record.}
\begin{enumerate}[resume]
\item What did Seasholes and Wu find after upper-limit events?
\item What did Chen and co-authors find about large investors?
\item What did Jones and co-authors find about retail accounts of different sizes?
\item What did Bian and co-authors find about \omterm{def:s1:asia-specific-equity-strategies:northbound}{northbound flows}?
\end{enumerate}

\medskip\noindent\textbf{Part III --- The events.}
\begin{enumerate}[resume]
\item How often do limit-up closes happen at 5\%, 10\% and 20\% limits?
\item Give the overnight gain for every close and weighted by fills at a 10\% limit.
\item Why do fills cut the gain?
\item What does the magnet do to the distribution of daily moves?
\end{enumerate}

\medskip\noindent\textbf{Part IV --- The verdict.}
\begin{enumerate}[resume]
\item State the \emph{named result}: the next-day return after a limit-up lock and how it changes with the limit's width.
\item Decompose the overnight gain into backlog and attention.
\item Why is the retail-attention reversal hard to harvest in the mainland?
\item What did the 2024 data change do to the northbound-flow signal?
\item How would you backtest a limit-up strategy honestly?
\item Which strategy file depends most on regulation?
\item How does this chapter relate to chapter~9 on flows?
\item In one sentence: what does a price limit sell, and to whom?
\end{enumerate}
\end{problem}

\section{Interview questions}

\begin{interviewq}[$\star$ \iqroles{researcher}]\label{iq:s1:asia-specific-equity-strategies:1}
Why is a momentum signal fitted on the closing prices of a market with daily limits suspect?
\end{interviewq}

\begin{interviewq}[$\star\star$ \iqroles{trader}]\label{iq:s1:asia-specific-equity-strategies:2}
A stock you want to buy is locked limit up with a large queue. What do you do?
\end{interviewq}

\begin{interviewq}[$\star\star$ \iqroles{researcher}]\label{iq:s1:asia-specific-equity-strategies:3}
How would you detect the \omterm{def:s1:asia-specific-equity-strategies:magnet}{magnet effect} in intraday data?
\end{interviewq}

\begin{interviewq}[$\star\star$ \iqroles{developer}]\label{iq:s1:asia-specific-equity-strategies:4}
What changes in a backtester for a market with daily limits and the T+1 rule?
\end{interviewq}

\begin{interviewq}[$\star\star$ \iqroles{risk}]\label{iq:s1:asia-specific-equity-strategies:5}
What are the risks of holding a book of mainland shares that may lock limit down?
\end{interviewq}

\begin{interviewq}[$\star\star\star$ \iqroles{researcher}]\label{iq:s1:asia-specific-equity-strategies:6}
A stock's value jumps by $x$ above a limit $L$ (in log terms). Buy orders queue at the limit and fill with probability $e^{-(x - L)/s}$ for $x > L$;
the next open reflects the full value. Derive the expected overnight gain of a filled order for $x - L$ exponentially distributed with mean $m$, and
compare it with the unconditional mean.
\end{interviewq}
